% ============================================================================
% Fichier  : partie2/P02_C09_facteur_humain.tex
% Titre    : Le Facteur Humain
% Auteur   : MINKA MI NGUIDJOÏ Thierry Emmanuel
% Version  : 1.0 -- novembre 2026
% ============================================================================
% Base : brouillon « Livre II » (rôles) et « Livre VIII » (sensibilisation).
% Sources : ISO/IEC 27002:2022 (mesures 6.1 à 6.8) ; Cialdini (1984) ;
% Cappelli, Moore et Trzeciak (2012) ; NIST SP 800-50 (2003).
% ============================================================================

\chapter{Le Facteur Humain}
\label{chap:humain}
\chaptermark{Facteur humain}

\epigraph{%
    « L'utilisateur n'est pas le maillon faible~:\\
    c'est le maillon que l'on a oublié de former. »%
}{Principe du cours}

\niveauChapitre{Fondamental}{Chapitres~\ref{chap:menaces}
et~\ref{chap:gouvernance}.}

\begin{boxobjectifs}
\begin{itemize}
    \item Expliquer pourquoi l'ingénierie sociale fonctionne, à partir
          des leviers de l'influence.
    \item Organiser la sécurité avant, pendant et après l'emploi.
    \item Concevoir un programme de sensibilisation et en mesurer l'effet.
    \item Conduire une campagne d'hameçonnage simulé dans un cadre éthique.
    \item Prévenir et détecter la menace interne.
\end{itemize}
\end{boxobjectifs}

\bigskip

% ============================================================================
\section*{Question directrice : le technicien de Garoua était-il fautif~?}
% ============================================================================

Le technicien de SAHEL INFOGÉRANCE a donné son mot de passe à une fausse
page. Le blâmer est tentant, et inutile~: personne ne l'avait formé, aucun
second facteur ne protégeait son compte, et il recevait chaque jour des
dizaines de courriels urgents. La sécurité qui repose sur la vigilance
parfaite de chacun échoue toujours~; celle qui combine formation, mesures
techniques et culture du signalement tient.

\medskip
\textbf{Plan du chapitre.}
\begin{enumerate}
    \item Pourquoi l'ingénierie sociale fonctionne (section~\ref{sec:c09-influence}).
    \item La sécurité au fil de l'emploi (section~\ref{sec:c09-rh}).
    \item Sensibiliser et former (section~\ref{sec:c09-sensibiliser}).
    \item L'hameçonnage simulé (section~\ref{sec:c09-simulation}).
    \item La menace interne (section~\ref{sec:c09-interne}).
\end{enumerate}

\separateur

% ============================================================================
\section{Pourquoi l'Ingénierie Sociale Fonctionne}
\label{sec:c09-influence}
% ============================================================================

R.~Cialdini a décrit les leviers qui font qu'une personne accepte une
demande sans l'examiner~\cite{cialdini1984}. Les attaquants les
utilisent tous.

\begin{tableaucours}{L{2.8cm}L{5.0cm}L{4.4cm}}{%
    \tableauHeader{Levier} &
    \tableauHeader{Principe} &
    \tableauHeader{Usage par l'attaquant}
}
Autorité & On obéit à qui semble légitime & Faux directeur, faux
technicien \\
Urgence et rareté & On agit vite quand le temps manque & « Votre compte
sera suspendu dans 24~h » \\
Réciprocité & On rend ce que l'on reçoit & Un « cadeau », puis une
demande \\
Sympathie & On aide qui nous plaît & Faux collègue chaleureux \\
Preuve sociale & On fait comme les autres & « Tous vos collègues ont déjà
validé » \\
Engagement & On reste cohérent avec ce qu'on a commencé & Une petite
demande, puis une plus grande \\
\end{tableaucours}
\vspace{-0.8em}
\begin{center}{\small\itshape Tableau~9.1. Les leviers de l'influence et
l'ingénierie sociale}\end{center}

% ============================================================================
\section{La Sécurité au Fil de l'Emploi}
\label{sec:c09-rh}
% ============================================================================

L'ISO/IEC~27002 regroupe les mesures liées aux personnes autour de trois
moments~\cite{iso27002_2022}.

\begin{tableaucours}{L{2.4cm}L{9.8cm}}{%
    \tableauHeader{Moment} &
    \tableauHeader{Mesures}
}
Avant & Vérification des antécédents proportionnée au poste~; clauses de
confidentialité et de sécurité dans le contrat \\
Pendant & Information sur les règles~; formation~; procédure
disciplinaire connue~; signalement des incidents \\
Changement de poste & Revue des droits~: on retire les anciens avant
d'ajouter les nouveaux \\
Départ & Retrait de tous les accès le jour du départ~; restitution du
matériel~; rappel des obligations qui survivent au contrat \\
\end{tableaucours}
\vspace{-0.8em}
\begin{center}{\small\itshape Tableau~9.2. La sécurité des ressources humaines}\end{center}

\begin{boxalert}
\textbf{Violation typique.} Les comptes d'un salarié parti restent actifs
pendant des mois.

\smallskip
\textbf{Règle~:} les ressources humaines informent l'informatique de
chaque départ, au plus tard le jour même~; une revue mensuelle rapproche
la liste des comptes de celle du personnel (Chapitre~\ref{chap:iam}).
\end{boxalert}

% ============================================================================
\section{Sensibiliser et Former}
\label{sec:c09-sensibiliser}
% ============================================================================

Le NIST distingue la \emph{sensibilisation} (faire connaître les risques
à tous), la \emph{formation} (donner des compétences à ceux qui en ont
besoin) et l'\emph{éducation} (former des spécialistes)~\cite{nist80050}.

\begin{boxprocedure}
\textbf{Construire un programme de sensibilisation}
\begin{enumerate}
    \item \textbf{Partir des risques}~: hameçonnage, mots de passe, fraude
          au président, supports amovibles, signalement.
    \item \textbf{Cibler}~: un module commun, et des modules par population
          (comptabilité, administrateurs, prestataires).
    \item \textbf{Faire court et répété}~: des séquences de dix minutes
          plusieurs fois par an valent mieux qu'une journée annuelle.
    \item \textbf{Donner un geste simple}~: « en cas de doute, je signale à
          tel numéro ».
    \item \textbf{Mesurer}~: taux de participation, taux de clic et taux
          de signalement lors des simulations.
    \item \textbf{Récompenser le signalement}, ne jamais le punir.
\end{enumerate}
\end{boxprocedure}

% ============================================================================
\section{L'Hameçonnage Simulé}
\label{sec:c09-simulation}
% ============================================================================

Une campagne d'hameçonnage simulé envoie de faux messages au personnel
pour mesurer et améliorer ses réflexes. Elle se conduit dans un cadre
écrit~: approbation de la direction, information préalable des
représentants du personnel, aucune sanction individuelle, résultats
agrégés, et une page pédagogique pour qui a cliqué.

\begin{boxscenario}{Opération WOURI. Trois campagnes à la BANQUE SANAGA}
\textbf{Constatations} (résultats agrégés de trois campagnes sur
200~agents)~:

\medskip
\begin{tableaucours}{L{3.0cm}ccc}{%
    \tableauHeader{Campagne} &
    \tableauHeader{Ont cliqué} &
    \tableauHeader{Ont saisi un mot de passe} &
    \tableauHeader{Ont signalé}
}
Janvier & 62 & 31 & 8 \\
Avril & 38 & 15 & 41 \\
Juillet & 24 & 7 & 77 \\
\end{tableaucours}

\medskip
\textbf{Analyse} (confiance moyenne). Le taux de clic passe de 31~\% à
12~\%, et surtout le taux de signalement de 4~\% à 38,5~\%. Le second
indicateur est le plus utile~: un signalement rapide permet de bloquer une
vraie attaque avant qu'elle ne s'étende. La confiance est moyenne car les
messages simulés différaient d'une campagne à l'autre.
\end{boxscenario}

% ============================================================================
\section{La Menace Interne}
\label{sec:c09-interne}
% ============================================================================

\begin{boxdef}
\termecle{Menace interne}~: risque qu'une personne disposant d'un accès
légitime (salarié, prestataire, ancien salarié) l'utilise, volontairement
ou non, au détriment de l'organisation.
\end{boxdef}

Les travaux du CERT de l'université Carnegie Mellon, fondés sur des
centaines de cas, distinguent le sabotage, le vol de propriété
intellectuelle et la fraude, et montrent que des signaux précèdent souvent
l'acte~: conflit avec la hiérarchie, accès hors des horaires habituels,
téléchargements inhabituels avant un départ~\cite{cappelli2012}.

\begin{tableaucours}{L{4.0cm}L{8.2cm}}{%
    \tableauHeader{Mesure} &
    \tableauHeader{Effet}
}
Moindre privilège & Limite ce qu'un initié peut faire \\
Séparation des tâches & Empêche qu'une seule personne commette et cache
une fraude \\
Journalisation des accès sensibles & Permet de détecter et de prouver \\
Revue des droits & Retire les accès devenus inutiles \\
Procédure de départ & Ferme les accès avant qu'ils ne servent \\
\end{tableaucours}
\vspace{-0.8em}
\begin{center}{\small\itshape Tableau~9.3. Prévenir la menace interne}\end{center}

\begin{boxmetier}{RSSI}
\textbf{Une culture, pas une police.} Le RSSI fait du signalement un
réflexe valorisé. Il mesure le délai entre un message suspect et son
signalement, et le présente au comité de sécurité comme un indicateur de
santé.
\end{boxmetier}

\begin{boxmetier}{Analyste SOC}
\textbf{Traiter chaque signalement.} L'analyste répond à chaque
signalement, même faux~: un remerciement rapide entretient le réflexe. Un
vrai message d'hameçonnage signalé déclenche la recherche des autres
destinataires et le blocage du domaine.
\end{boxmetier}

\separateur

\begin{boxresume}
\textbf{Ce qu'il faut retenir de ce chapitre~:}
\begin{itemize}
    \item \textbf{L'ingénierie sociale} exploite des leviers humains
          normaux~: autorité, urgence, réciprocité, sympathie, preuve
          sociale, engagement.
    \item \textbf{La sécurité des personnes} couvre l'avant, le pendant et
          l'après emploi~; les accès d'un salarié parti se ferment le jour
          même.
    \item \textbf{La sensibilisation} est courte, répétée, ciblée et
          mesurée~; le signalement se récompense.
    \item \textbf{L'hameçonnage simulé} se conduit dans un cadre écrit, sans
          sanction individuelle~; le taux de signalement est l'indicateur
          clé.
    \item \textbf{La menace interne} se prévient par le moindre privilège,
          la séparation des tâches, la journalisation et les revues.
\end{itemize}
\end{boxresume}

\begin{boxinfo}
\textbf{Points de contrôle de l'élève.} Je sais~: nommer les six leviers
de l'influence~; organiser la sécurité au départ d'un salarié~; concevoir
un programme de sensibilisation~; fixer le cadre d'une campagne
d'hameçonnage simulé~; citer trois signaux de menace interne.
\end{boxinfo}

\bigskip

\begin{boxexo}[title={\ding{45}\quad Exercice~9.1. Les leviers \niveaubase}]
Identifiez les leviers de l'influence dans le courriel du 19~février
(Chapitre~\ref{chap:menaces}).
\end{boxexo}

\begin{boxexo}[title={\ding{45}\quad Exercice~9.2. Le départ \niveaubase}]
Listez, dans l'ordre, les actions à mener le jour du départ d'un
administrateur systèmes.
\end{boxexo}

\begin{boxexo}[title={\ding{45}\quad Exercice~9.3. Les taux \niveaubase}]
Recalculez, à partir du tableau de l'Opération WOURI, les taux de clic,
de saisie et de signalement des trois campagnes.
\end{boxexo}

\begin{boxexo}[title={\ding{45}\quad Exercice~9.4. Changement de poste \niveaumoyen}]
Mme~OWONA passe de la saisie des paiements à leur validation. Que doit-il
se passer pour ses droits, et pourquoi~?
\end{boxexo}

\begin{boxexo}[title={\ding{45}\quad Exercice~9.5. Le cadre \niveaumoyen}]
Rédigez la note qui autorise une campagne d'hameçonnage simulé à l'ANPP~:
objectifs, population, règles, traitement des résultats.
\end{boxexo}

\begin{boxexo}[title={\ding{45}\quad Exercice~9.6. Programme \niveaumoyen}]
Proposez un programme annuel de sensibilisation pour LiZbiZ~: thèmes,
populations, formats, calendrier, indicateurs.
\end{boxexo}

\begin{boxexo}[title={\ding{45}\quad Exercice~9.7. Signaux \niveaumoyen}]
Un comptable de LiZbiZ, en conflit avec sa direction, télécharge en une
nuit l'historique des paiements de cinq ans, deux semaines avant sa
démission. Quels signaux, quelles mesures auraient permis de le voir~?
\end{boxexo}

\begin{boxexo}[title={\ding{45}\quad Exercice~9.8. Sanctionner~? \niveaumoyen}]
Le directeur de la BANQUE SANAGA veut sanctionner les agents qui ont
cliqué. Que lui répondez-vous~?
\end{boxexo}

\begin{boxexo}[title={\ding{45}\quad Exercice~9.9. Le prestataire \niveauexpert}]
Les prestataires sont rarement formés par le client. Proposez les clauses
et les mesures qui étendent le programme humain de LiZbiZ à SAHEL
INFOGÉRANCE.
\end{boxexo}

\begin{boxexo}[title={\ding{45}\quad Exercice~9.10. Fraude au président \niveauexpert}]
Rédigez la procédure qui protège la direction financière de LiZbiZ
contre la fraude au président~: vérifications, seuils, double validation,
contre-appel.
\end{boxexo}

\bigskip

\begin{boxinfo}
\textbf{Pour aller plus loin.}
\begin{itemize}
    \item R.~Cialdini, \emph{Influence}~\cite{cialdini1984}.
    \item D.~Cappelli, A.~Moore et R.~Trzeciak, \emph{The CERT Guide to
          Insider Threats}~\cite{cappelli2012}.
    \item NIST SP~800-50, programme de sensibilisation~\cite{nist80050}.
\end{itemize}
\end{boxinfo}

\begin{boxinfo}
\textbf{Transition.} Gouvernance, risques, actifs et personnes forment le
socle organisationnel. La Partie~\ref{part:normatif} présente les normes
qui le structurent et le droit qui l'encadre.
\end{boxinfo}
